\section{Branch-count and quadratic-field refinements of standard descent}
\label{sec:standard-partner-hurwitz-refinement}

\paragraph{Author freeze and provenance.}
HD-SS08--10, reserved question Q2eh, is a new derivation by
\texttt{/root/sixfolds} from the accepted HD-SS01--07.
Root requested a separate check of the surviving positive-partner
case after HD-SS06--07. The branch-count refinement and the exact
quadratic-field parity argument are written in full below.
The complete author, root and independent whole-proof readings pass,
including the source interfaces and author checker/first output.
Numerical consequences are distinguished
from the structural description of the surviving case. No elimination
of the remaining unbounded partner degrees is asserted.

Use the actual sphere setting of HD-SS01--02:
\[
 V=(\operatorname{Std}_n\circ x)\otimes U,\qquad
 x(J)=P\in\{A_n,S_n\},\qquad b=\dim U\ge2,
 \qquad \mu_V=6.                                      \tag{SH.1}
\]
The representation V is faithful and absolutely irreducible, both
factors are actual finite linear representations, and the original
cover is a connected J-Galois cover \(C\to\mathbf P^1\).
Put \(H=x^{-1}((S_n)_1)\), \(X=C/H\), \(h=g(X)\),
\(K_0=\ker U\), \(Q=U(J)\), and \(C_U=C/K_0\).
Write \(\mu=\mu_U\). When \(K_0\) is transitive, put
\(N=H\cap K_0\) and \(Y=C/N\).
The exact quadratic model, period-rank and natural-polarization
conditions will be imposed only where stated.

\subsection*{HD-SS08. The complete positive-partner branch table}

\begin{lemma}[HD-SS08]
Assume \(n\ge5\), \(K_0\) is transitive and \(\mu>0\).
Let \(s\) be the genuine branch count of the actual
Q-cover \(C_U\to\mathbf P^1\), with sorted inertia orders
\(d_1,\ldots,d_s\), and put
\[
 \Delta=-2+\sum_{i=1}^s(1-1/d_i),\qquad
 f_i=\dim\ker(U(q_i)-1).
 \tag{SH.2}
\]
Then \(\Delta>0\), \(s\ge3\), and the following bounds hold.
The column for \(h\ge0\) includes every possible natural genus.
\[
\begin{array}{c|c|c|c|c}
\text{actual signature case}&\Delta\text{ lower bound}
 &\mu\text{ upper bound}&n\ (h\ge0)&n\ (h\ge1)\\ \hline
(2,3,7)&1/42&b-1&126(b+1)&42(b+5)\\
s=3,\ d_1=2,\ (d_i)\ne(2,3,7)&1/24&b-1&72(b+1)&24(b+5)\\
s=3,\ d_1\ge3&1/12&b&36b+48&12b+72\\
s=4&1/6&2b&24b+24&12b+36\\
s\ge5&(s-4)/2&(s-2)b&10b+8&6b+12
\end{array}                                             \tag{SH.3}
\]
In particular \(n\le126(b+1)\), and \(h\ge1\) gives
\(n\le42(b+5)\). Neither an exact coefficient-field assumption
nor period rank nine is needed for this lemma.
\end{lemma}

\begin{proof}
HD-SS01 proves U irreducible. Transitivity gives \(J=K_0H\),
\(H/N\simeq Q\), and the actual finite map \(Y\to C_U\)
has degree n. Positive \(\mu\) puts U in \(H^1(C_U,\mathbf C)\).
Genus zero has no such cohomology; genus one has only its invariant
holomorphic and antiholomorphic lines, whose irreducible constituents
have degree one. Thus \(g(C_U)\ge2\) since \(b\ge2\).
The orbifold excess is therefore positive.

Euler characteristic of the actual Q-sphere cover gives
\[
 \mu=(s-2)b-\sum_{i=1}^s f_i.
 \tag{SH.4}
\]
There is no invariant-vector correction because U is nontrivial
irreducible. If \(s=3\), no inertia matrix is scalar. Indeed, if
one were scalar, the product-one relation would make the other two
commute; all three would commute, and the group they generate would
have only one-dimensional complex irreducible representations.
This contradicts \(b\ge2\). In particular an order-two inertia
has both eigenvalues 1 and -1, hence \(f_i\ge1\).
This proves the \(b-1\) upper bound in the first two rows.
All other upper bounds for \(\mu\) follow by dropping the
nonnegative fixed-space dimensions in (SH.4).

For clarity the entire excess calculation is elementary.
At three sorted orders, \(d_1=d_2=2\) is never hyperbolic.
If \((d_1,d_2)=(2,3)\), then \(d_3\ge7\); order seven
gives \(1/42\), and order at least eight gives at least \(1/24\).
If \((d_1,d_2)=(2,4)\), then \(d_3\ge5\), giving at least
\(1/20\). If \(d_1=2,d_2\ge5\), the excess is at least
\(1/10\). With no order two, \((3,3,3)\) has zero excess,
and the least positive case is \((3,3,4)\), of excess \(1/12\).
Larger sorted orders give at least that value. At four marks all
orders two give zero, and otherwise replacing one order by three
gives the least positive value \(1/6\). At least five marks give
\(\Delta\ge(s-4)/2\).

The actual quotient and Riemann--Hurwitz calculation of (SSK.8),
or directly the two maps \(Y\to C_U\) and \(Y\to X\), gives
\[
 n\Delta\le 6+\mu-2(b-1)(h-1).
 \tag{SH.5}
\]
Every genuine Q-inertia on Y has positive moving rank under its
faithful representation, which bounds its orbifold contribution.
Equation (SS.9) supplies the last Euler equality in this calculation.
For \(h\ge0\), the right side is at most \(\mu+2b+4\);
for \(h\ge1\), it is at most \(\mu+6\).
Substitute the preceding bounds to obtain the first four rows.
For the last row the exact bounds are
\[
 n\le \frac{2sb+8}{s-4}\le10b+8\quad(h\ge0),\qquad
 n\le \frac{2(s-2)b+12}{s-4}\le6b+12\quad(h\ge1).
 \tag{SH.6}
\]
Both final inequalities follow by writing the respective quotients
as \(2b+(8b+8)/(s-4)\) and \(2b+(4b+12)/(s-4)\).
All entries of (SH.3) now follow, including all possible branch counts.
\end{proof}

\subsection*{HD-SS09. Exact quadratic models and improved cutoffs}

Now suppose V has a model over its exact imaginary quadratic
character field K. Retain the transitivity hypothesis when applying
the next local calculation. HD-SS01 then proves that U and
\(\tau=U|_H\) have models over the same exact field K.
It is essential that these are actual K-models; a model only over
a larger field is not substituted.

\begin{proposition}[HD-SS09]
Under the positive-partner and transitivity hypotheses of HD-SS08,
with the exact K-model just specified,
\[
                     n\le126b+84,\qquad
             h\ge1\ \Longrightarrow\ n\le42b+168.
 \tag{SH.7}
\]
In the actual \((2,3,7)\) case there is the sharper residue bound
\[
 \mu\le b-\varepsilon_K(b),\qquad
 \varepsilon_K(b)=
 \begin{cases}
  \eta(b)+(b\bmod6),&K=\mathbf Q(\sqrt{-3}),\\
  2+(b\bmod3),&K=\mathbf Q(\sqrt{-7}),\\
  2+(b\bmod6),&K\ne\mathbf Q(\sqrt{-3}),\mathbf Q(\sqrt{-7}),
 \end{cases}
 \quad
 \eta(b)=\begin{cases}1,&b\text{ odd},\\2,&b\text{ even}.
 \end{cases}
 \tag{SH.8}
\]
Remainders in (SH.8) are the least nonnegative ones; in every case
\(\varepsilon_K(b)\ge2\).

Consequently, for actual rank-nine sources in this exact-K setting,
without a separate kernel-transitivity hypothesis,
\[
 h\ge1\ \Longrightarrow\ n\le42b+168.
 \tag{SH.9}
\]
For a source of smallest coefficient degree in the exact comparison
class of HD-SS06--07, one also has
\[
                              n\le126b+84.
 \tag{SH.10}
\]
The comparison class consists of actual exact-K, rank-nine sources
with the specified natural signed polarization determinant. The
last conclusion is not a bound on all nonminimal realizations.
\end{proposition}

\begin{proof}
In the \((2,3,7)\) case Q is a quotient of the corresponding
triangle group and is perfect. In its abelianization write the
generators additively as u,v,w, with \(2u=3v=7w=0\) and
\(u+v+w=0\). Then \(2v=-2w\) has order dividing both three
and seven, so it is zero. Hence v,w,u all vanish. Every quotient
abelianization is therefore trivial.

The determinant character of U is trivial. Write \(f_2,f_3,f_7\)
for the three fixed-space dimensions. The nonscalar involution has
a positive fixed space, and its -1 multiplicity is even. Thus
\[
                         f_2\ge\eta(b),\qquad f_2\equiv b\pmod2.
 \tag{SH.11}
\]
Because the characteristic polynomials are over K, root multiplicities
are constant on Galois orbits over K. Unless \(K=\mathbf Q(\sqrt{-3})\),
the two nontrivial cube roots form one orbit, giving
\(f_3\ge b\bmod2\). For seventh roots the orbit length is six
unless \(K=\mathbf Q(\sqrt{-7})\), in which case it is three.
Therefore \(f_7\ge b\bmod6\) or \(f_7\ge b\bmod3\), respectively.

These elementary field assertions require no character classification.
For completeness, \(\Phi_7(z+1)\) is Eisenstein at seven, so the
cyclotomic field has degree six, with automorphisms indexed by the
cyclic group \(\mathbf F_7^\times\). Its unique quadratic subfield
is \(\mathbf Q(\sqrt{-7})\): if
\(a=\zeta+\zeta^2+\zeta^4\) and
\(\bar a=\zeta^3+\zeta^5+\zeta^6\), then
\(a+\bar a=-1\), \(a\bar a=2\), and \((a-\bar a)^2=-7\).
The two cubic factors are
\((z^3-az^2+\bar a z-1)\) and its conjugate.
A different quadratic K has trivial intersection with this
Galois field, so the nontrivial seventh roots still form a six-element
orbit over K. The cube-root assertion follows directly from
\(\Phi_3=z^2+z+1\).

If K is Eisenstein, use \(f_3\ge0\) and the other two lower bounds.
If K is Gaussian or any quadratic field distinct from the two
exceptions, use \(f_2+f_3\ge2\) and the six-orbit bound for
\(f_7\). For \(K=\mathbf Q(\sqrt{-7})\), use the same
\(f_2+f_3\ge2\) and the three-orbit bound. Since
\(\mu=b-f_2-f_3-f_7\), these are exactly (SH.8).
For Eisenstein K, even b gives \(f_2\ge2\), and odd b gives
\(f_2\ge1\) and \(b\bmod6\ge1\). Thus all three bounds
indeed have \(\varepsilon_K(b)\ge2\).

Inserting \(\mu\le b-2\) into (SH.5), with
\(\Delta=1/42\), proves (SH.7) in the triangle case.
Every other row of (SH.3) is smaller: for \(h\ge0\) its
largest bound is \(72(b+1)<126b+84\), and for \(h\ge1\)
it is at most \(24(b+5)<42b+168\). This proves (SH.7).

To remove the transitivity premise in (SH.9), suppose its claimed
bound fails. Then \(n>42b+168>\max(b+1,72)\), so HD-SS07
forces transitive \(K_0\). Actual period rank nine and \(h\ge1\)
give \(\mu\ge3\) by (SS.16). The positive-partner bound (SH.7)
is a contradiction.

Finally suppose a minimal source violates (SH.10). HD-SS07 again
gives transitivity. Positive \(\mu\) contradicts (SH.7); hence
\(\mu=0\). The post-vanishing portion of HD-SS06 now applies:
exact K, the rational base, the primitive-line correspondence,
positive rational polarization scale and actual rank-nine descent
use only transitivity and \(\mu=0\), not the earlier constant 168.
It constructs an actual source of coefficient degree b instead of
\((n-1)b\), with the same natural signed determinant, contradicting
the stated minimality. This proves (SH.10).
\end{proof}

\subsection*{HD-SS10. The remaining high-ratio source is a triangle pullback}

\begin{theorem}[HD-SS10]
In the exact-K rank-nine setting above, assume the source is minimal
in the HD-SS06 comparison class and
\[
                              n>72(b+1).
 \tag{SH.12}
\]
Then \(\mu>0\), the actual Q-cover \(C_U\to\mathbf P^1\)
has exactly the signature \((2,3,7)\), and Q is perfect.
Moreover, under the faithful map \(g\mapsto(x(g),U(g))\),
\[
 J=P\times Q,\qquad h=0,\qquad b\ge7,
 \qquad
 \max\{1,18-2b\}\le\mu\le b-\varepsilon_K(b).
 \tag{SH.13}
\]
The connected cover \(Y\to X=\mathbf P^1\) is the normalized
pullback of the rigid three-point Q-cover along the actual degree-n
rational map \(X\to\mathbf P^1\). In particular the necessary
remaining numerical interval is
\[
              72(b+1)<n\le42\bigl(3b+4-\varepsilon_K(b)\bigr).
 \tag{SH.14}
\]
These are necessary conditions on a residual class. They neither
construct such a source nor prove that its V-period rank is smaller
than nine.
\end{theorem}

\begin{proof}
HD-SS07 makes the kernel transitive in (SH.12). If \(\mu=0\),
the same actual polarized descent used above contradicts minimality.
Thus \(\mu>0\). All rows of (SH.3) other than \((2,3,7)\)
contradict (SH.12), so that signature is forced. Perfectness and
the exact quadratic-field bounds have already been proved.

The map \(J\to P\times Q\) is injective by faithfulness of V.
The transitive-kernel argument of HD-SS06 supplies
\(A_n\times1\subset J\). If \(P=A_n\), surjectivity to Q
immediately gives the full direct product. If \(P=S_n\), the
image modulo \(A_n\times1\) is a subdirect product of
\(C_2\times Q\). Its kernel over Q is either all of \(C_2\)
or trivial. In the latter case it is the graph of a surjective
homomorphism \(Q\to C_2\), impossible because Q is perfect.
Thus the image is the full product, and so is J. This is an
actual group identity, not merely an abstract quotient assertion.

If \(h\ge1\) and \(b\ge4\), (SH.9) contradicts (SH.12),
since \(72(b+1)>42b+168\). If \(b=2\) or three, the required
\(\mu\ge3\) from (SS.16) contradicts \(\mu\le b-2\).
Hence \(h=0\).

The descended representation \(\tau\) has image Q, degree b and
trivial determinant. Every genuine inertia matrix in it has
moving rank at least two: a finite-order nonidentity matrix of
moving rank one would have a single eigenvalue different from one, equal
to its determinant. The V-variation is a fixed direct summand
of the tau-variation by (SS.7). Its actual rank nine forces at
least twelve genuine marks on the rational base X. Consequently
\[
                    24\le D'=2b+6+\mu\le3b+4,
 \tag{SH.15}
\]
which proves \(b\ge7\) and the lower bound in (SH.13).

In the actual direct product, \(H=P_1\times Q\),
\(K_0=P\times1\), and \(N=P_1\times1\). The fixed fields
of H and \(K_0\) have compositum the fixed field of N.
This identifies Y with the normalization of their fiber product
over the original sphere, proving the pullback assertion.
The three marked Q-cover is fixed after ordering and normalizing
its marks and making a finite base change, by the finite Hurwitz
framework for a fixed finite group and three branch points.
This rigidity refers only to the partner quotient. It does not
make the whole pullback family or the V summand constant.
Finally (SH.5), with \(h=0\), \(\Delta=1/42\) and (SH.8),
gives (SH.14).
\end{proof}

\paragraph{Source and comparison boundary.}
All curve, Hodge and rational-polarization inputs are the accepted
HD-SS01--07 and D-NR, with their original proof/source status.
The positive-excess case table, triangle abelianization, cyclotomic
orbit calculation and group-product argument are proved here.
No new generic-endomorphism, simple-monodromy, classification or
unrefereed-preprint premise is used. The numerical cutoffs refine
the already accepted standard descent; the structural theorem
describes the positive-partner residue but leaves its large-b
existence and period-rank question open.
